\section{Introduction}

Recent advancements in video generation have led to the development of methods aimed at synthesizing long-form videos with rich narrative storylines. 
Autoregressive techniques \cite{guo2025long,henschel2024streamingt2v,qiu2023freenoise,song2025dfot,causvid} have made significant progress in extending temporal horizons, yet they still face challenges in producing complex, story-driven videos. 
While these methods improve consistency in generating adjacent frames, they are often limited to single-shot or repetitive-scene settings, which leads to repetitive content and shallow narratives. 
Furthermore, autoregressive approaches struggle with error accumulation, weak global planning, and insufficient long-range coherence, hindering their ability to generate videos with evolving storylines and seamless scene transitions.

An alternative and more effective approach is story-frame generation, where story frames are first created using methods similar to text-to-image techniques, and then image-to-video models are used to interpolate intermediate frames to construct the full video. This strategy offers a more controlled and flexible way to generate story-driven content, allowing for better narrative progression across multiple scenes. 
However, existing story-frame generation methods face significant limitations. IC-LoRA\cite{huang2024context}, for example, is highly constrained by its ability to generate only limited frames at a time, severely limiting its flexibility and ultimately degrading the quality of the story frames. Meanwhile, methods like StoryDiffusion\cite{zhou2024storydiffusion}, VideoGen-of-Thought\cite{zheng2025videogenofthoughtstepbystepgeneratingmultishot}, and One Prompt One Story\cite{liu2025one}, while more flexible, heavily focus on character-centric shots, leading to repetitive perspectives and weak story development. Additionally, as the number of story frames increases, challenges like consistency, semantic alignment, and narrative coherence become more pronounced, further diminishing the overall video quality.

Recent multimodal large language models (MLLMs) like GPT-4o\cite{gpt4o} can generate multiple story frames but require detailed prompts and often struggle with consistency and producing the requested number of frames. Over time, reliance on historical context leads to accumulating inconsistencies and a lack of global narrative coherence, hindering the creation of high-quality, multi-scene story-driven frames.


\begin{figure}[t]
    \centering
    \includegraphics[width=0.9\textwidth]{fig/intro_new.png} 
    \caption{The story frames generated by \ours{} exhibit high narrative quality and subject consistency across multiple scenes.}
    \label{fig:intro}
\end{figure}

To address the limitations of existing methods for generating high-narrative, multi-scene, and story-driven frames under strong consistency, we propose \ours{}, a unified framework that captures key narrative events across evolving scenes while maintaining temporal coherence. Built on the Step-Video-T2V \cite{ma2025stepvideot2vtechnicalreportpractice,huang2025step} backbone, \ours{} decouples story planning from frame synthesis, generating editable visual anchors. A bidirectional story frame predictor ensures long-range dependencies like character continuity and scene transitions are maintained, as shown in \Cref{fig:intro}. We introduce a Multi-Event Story Frame Labeling, which generates temporally structured, semantically rich captions that reflect both global story structure and detailed event-level dynamics. Progressive Story Frame Training, starting with global prompts and refining specific events, ensuring \ours{} generates high-quality, consistent, and flexible story-driven frames. Additionally, all story frames are designed to be editable and expandable, enabling the creation of longer sequences and allowing for manual narrative adjustments, significantly enhancing flexibility and user control.

Extensive experiments demonstrate that \ours{} outperforms existing story-frame generation models across multiple dimensions, including narrative quality, subject consistency, prompt-frame alignment, and diversity. Compared to open-source models, \ours{} achieves superior performance in these key areas. Furthermore, its results in narrative consistency and story richness are comparable to the closed-source GPT-4o, underscoring its ability to generate strong storytelling, coherent and compelling story frames.

Our contributions are summarized below:
\begin{itemize}
    \item We propose the \ours{} framework, which is built on a text-to-video model with controlling conditions to enable Bidirectional Story Frame Prediction, trained with diverse levels of event data from Multi-Event Story Frame Labeling pipeline, and a novel gradual strategy as Progressive Story Frame Training. 
    \item Our \ours{} can generate multiple high-consistency, multi-scene, and story-rich story frames in a single pass. Extensive experiments demonstrate the superior of \ours{} over existing open-source models considering narrative, consistency, alignment, and diversity of the generated frames. 
    \item \ours{} supports frame editing and expansion, providing a more flexible and controllable frame generation framework, allowing for longer sequences and manual narrative modifications.
   
\end{itemize}
